\section{树码的傅里叶结构}\label{sec:tree-fourier}

割依赖生成的词，因此每个保持该词的变换同样约束割的行傅里叶系数。
在每个门处，这些约束给出所有子节点对共享的一个公共傅里叶指标。
把门表取反即可区分该指标是否为求值标记。随后我们证明：在随机切片
处，把整个公共指标只保留在每对的一个子节点中，平滑后的傅里叶质量
损失很小。最后我们把保留下来的指标表示成在其他行固定后仍然自由的
行坐标。这就是 \Cref{sec:extraction} 使用的局部形式。
本节通篇假设 $f$ 是 \eqref{eq:function-of-code} 意义下的码的函数。

\subsection{行傅里叶支撑}

对每个节点 $v$，递归地定义其\emph{生成子} $P(v)$（一个叶子集）：
叶子处 $P(v)=\{v\}$，内部节点处
$P(v)=P(v_{1,0})\cup P(v_{1,1})$。
即生成子总是沿第一对子节点的两侧；与切片不同，它不含随机选择。
把同一个 $h\in H$ 加到 $P(v)$ 中的所有行上，会把该子树在 $a$ 处的
输出乘以 $(-1)^{h(a)}$。这由 Equation~\eqref{eq:gate-rule} 归纳
可得：把一对的两个子节点乘以同一符号，所有门输入 $z_g$ 不变，而
前因子乘以该符号。

对行傅里叶指标 $C=(q_i)_{i=1}^k\in(H^*)^k$，令
\[
 \chi_C(R)=\prod_i\chi_{q_i}(R_i),\qquad
 |C|=|\{i:q_i\ne0\}|,\qquad
 q_v=\sum_{i\in P(v)}q_i.
\]
把 $v$ 的子节点 $(g,b)$ 的指标记为 $q_{v,gb}$。按定义有
$q_v=q_{v,10}+q_{v,11}$。在固定表处，行展开及其平滑为
\begin{equation}
 f(R,F)=\sum_C\widehat f_C(F)\chi_C(R),\qquad
 \widehat{\bar f}_C(F)=r^{|C|}\widehat f_C(F).
 \label{eq:row-smoothing}
\end{equation}

\begin{lemma}[行对称性]
\label{lem:row-symmetry}
设 $f$ 是码的函数。$f$ 的每个非零行傅里叶系数在每个内部节点 $v$ 处
满足
\begin{equation}
 q_{v,g0}+q_{v,g1}=q_v\qquad(g\in[s]).
 \label{eq:row-relations}
\end{equation}
对表集合 $F$，记 $F^{v,-}$ 为仅把 $F_v$ 取反，$F^{v,z}$ 为把
$F_v(u)$ 换成 $F_v(u+z)$，其余表不变。则
\begin{align}
 \widehat f_C(F^{v,-})
   &=(-1)^{q_v(\one)}\widehat f_C(F),
   \label{eq:table-flip-action}\\
 \widehat f_C(F^{v,z})
   &=(-1)^{\sum_g z_g q_{v,g1}(\one)}\widehat f_C(F).
   \label{eq:table-translation-action}
\end{align}
以上陈述在行平滑后仍然成立。它们也被行傅里叶投影，以及到某个指定
整门表集合的表项总傅里叶次数范围的投影所保持。
\end{lemma}

\begin{proof}
对每个 $g$，把对 $g$ 中两个子节点的生成子都平移 $h$，对 $v$ 处子树
的效果与把 $P(v)$ 平移 $h$ 完全相同。因此两个行平移对整张词的效果
相同。取行傅里叶系数得
\[
 \chi_{q_{v,g0}+q_{v,g1}}(h)\widehat f_C(F)
   =\chi_{q_v}(h)\widehat f_C(F)\qquad(h\in H).
\]
不同特征是不同函数，证得 Equation~\eqref{eq:row-relations}。

把 $F_v$ 取反，其效果等同于把 $\one$ 加到 $P(v)$ 的所有行上。这证
明 Equation~\eqref{eq:table-flip-action}。为把 $F_v$ 的自变量平移
$z$，在每个 $z_g=1$ 的对中，把其第二个子节点的生成子加 $\one$。
这恰好改变那些门输入比特，而前因子不变。傅里叶变换给出
Equation~\eqref{eq:table-translation-action}。

行平滑与行傅里叶投影在 $C$ 上是对角的，故保持三个恒等式。在表项
符号的立方上，定义域平移置换坐标，而把门表取反则给每个单项式乘以
其奇偶符号。两种运算都保持任一指定整门表集合中的总次数，因此与
所述次数投影交换。
\end{proof}

\subsection{随机切片上的不分裂投影}

Equation~\eqref{eq:row-relations} 给每个子节点对指派同一和
$q_v$。Equation~\eqref{eq:table-flip-action} 由此有一个有用的解读：
在把 $F_v$ 取反下为奇的系数满足 $q_v(\one)=1$，故由
\eqref{eq:labels-dual} 其公共指标是一个标记。但单个子节点指标仍可以
是和为 $q_v$ 的任意泛函。

我们在一个门处化简这最后的自由度。把 $q_v$ 保留在每对的一个子节点
中，只剩下选哪个子节点的自由，而不是两个和给定的任意泛函。称 $C$
\emph{在 $v$ 处不分裂}，若对每个 $g$，有序对 $(q_{v,g0},q_{v,g1})$
为 $(q_v,0)$ 或 $(0,q_v)$。记 $S_v$ 为到这类指标的正交行傅里叶
投影。当 $q_v=0$ 时，这要求每个子节点指标为零。下一个界把代价对
切片深度取平均。它不要求指标在树的每个门处都不分裂。

\begin{lemma}[不分裂投影的平均代价]
\label{lem:unsplit}
设 $f$ 是满足 $|f|\le1$ 的码的函数，$r=1-\mu/d$（$0<\mu<1/2$），
$\bar f=T_rf$。选取与 $f$ 独立的随机切片。则
\begin{equation}
 \E_{J_0,\mathcal D}\sum_{v\in\mathcal D}
       \norm{(I-S_v)\bar f}_{2,R,F}^{2}
       \le \frac{s}{2e m\mu}.
 \label{eq:unsplit-bound}
\end{equation}
在决定 $f$ 的额外背景数据上取平均后，同样的界成立，只要切片的选取
独立于那些数据。
\end{lemma}

\begin{proof}
固定满足 Equation~\eqref{eq:row-relations} 的指标 $C$，令 $N_j$ 为
深度 $j$ 处指标 $q_v$ 非零的节点数。若 $q_v\ne0$，每个子节点对至少
含一个非零指标；该节点处不分裂条件的一次失败至少再提供一个非零
子节点。若 $q_v=0$，一次失败提供至少两个非零子节点，而其最小贡献
为零。因此深度 $j$ 处失败的个数 $B_j$ 满足
\[
 B_j\le N_{j+1}-sN_j.
\]
用切片包含概率与裂项相消得
\begin{align*}
 \E_{J_0,\mathcal D}\sum_{v\in\mathcal D}
          \one_{\{C\text{ 在 }v\text{ 处分裂}\}}
 &\le \frac1m\sum_{j=0}^{m-1}
          \frac{N_{j+1}-sN_j}{s^j}\\
 &=\frac{s}{m}\left(\frac{N_m}{s^m}-N_0\right)
 \le\frac{s|C|}{md}.
\end{align*}
此处 $N_m=|C|$。由 Equation~\eqref{eq:row-smoothing} 与 Parseval，
Equation~\eqref{eq:unsplit-bound} 的左边至多为
\[
 \frac{s}{m}\sum_C
   \left(r^{2|C|}\frac{|C|}{d}\right)
           \E_F|\widehat f_C(F)|^2.
\]
对 $x\ge0$ 有 $xe^{-2\mu x}\le(2e\mu)^{-1}$，且 $\log r\le-\mu/d$。
故每个括号因子至多 $(2e\mu)^{-1}$，而
$\sum_C\E_F|\widehat f_C(F)|^2=\norm f_2^2\le1$。
界得证。证明逐点依赖于背景数据，故其平均满足同一不等式。
\end{proof}

\subsection{自由平移与剩余行坐标}

抽取论证必须在固定大部分行之后使用公共指标。因此我们选取这样的
坐标：平移一个子节点的生成子只改变一个自由行，而其余坐标不动。
这样不分裂投影便在每个固定上下文处单独作用，且其保留的特征具有
显式形式。

固定切片 $\mathcal D$ 与节点 $v\in\mathcal D$。$2s$ 个集合
$P(v_{g,b})$ 非空且两两不交。按只依赖树的规则从每个集合中选一个
代表叶 $i_{g,b}$，令 $L_{g,b}=R_{i_{g,b}}$。对该生成子中的其他每个
叶子 $i$，把 $R_i$ 换成 $R_i+L_{g,b}$。这些生成子之外的行保持不变。
把所得的 $k-2s$ 个剩余坐标记为 $\Gamma\in H^{k-2s}$。

\begin{lemma}[自由平移坐标]
\label{lem:shift-coordinates}
坐标变换
\[
 H^k\longrightarrow H^{2s}\times H^{k-2s},\qquad
 R\longmapsto((L_{g,b}),\Gamma)
\]
是线性双射。特别地，在均匀行下，自由平移在给定 $\Gamma$ 时为独立
均匀。与完整行指标 $C$ 对应的 $L_{g,b}$ 的傅里叶指标是
$Q_{g,b}=q_{v,g b}$。

因此，在每个固定的 $\Gamma$ 处，算子 $S_v$ 在 $2s$ 个平移上的作用
恰是到不分裂平移指标的傅里叶投影。对码的函数或其行平滑，除非和
$Q_{g,0}+Q_{g,1}$ 对每个 $g$ 相同，否则平移系数为零。
Equations~\eqref{eq:table-flip-action} 与
\eqref{eq:table-translation-action} 的表作用在该固定上下文处对每个
平移系数成立，其中 $q_v=Q_{1,0}+Q_{1,1}$，$q_{v,g1}=Q_{g,1}$。
\end{lemma}

\begin{proof}
逆映射把代表行取为 $L_{g,b}$，并把其记录的差加 $L_{g,b}$ 恢复生成
子中的每一行。这证明双射性。在行特征中，$L_{g,b}$ 的系数恰是
$\sum_{i\in P(v_{g,b})}q_i=q_{v,g b}$。于是定义 $S_v$ 的条件只依赖
平移指标 $Q$，不依赖 $\Gamma$ 的任何傅里叶指标。把完整傅里叶展开
按 $Q$ 分组，即得其在每个固定纤维上的所述作用。同一分组证得支撑
与表作用恒等式：给定 $Q$ 的所有完整行系数满足同样的恒等式并带有
同样的符号。
\end{proof}

于是对非零公共指标 $q$，保留的平移指标由子集 $I\subseteq[s]$
（$q$ 出现在其中哪些对的第二个子节点）刻画。其特征为
\[
 \prod_{g\notin I}\chi_q(L_{g0})
 \prod_{g\in I}\chi_q(L_{g1}).
\]
若 $q=0$，只有全零的不分裂指标，其特征为 $1$。把 $F_v$ 取反会给
相应系数乘上 $(-1)^{q(\one)}$。在 $F_v$ 的奇部上，$q$ 是标记因而
非零，故由 $q$ 与 $I$ 的刻画唯一。对这样的标记，把 $F_v$ 的自变量
平移 $z\in\F^s$ 会给其系数乘上 $(-1)^{\sum_{g\in I}z_g}$。
因此同一子集 $I$ 同时记录行特征与表平移特征。这就是
\Cref{sec:extraction} 使用的局部傅里叶描述，在每个固定上下文处
有效。
